\documentclass[conference]{IEEEtran}

% --- Grundkonfiguration ---
\usepackage[utf8]{inputenc}
\usepackage[T1]{fontenc}
\usepackage[ngerman]{babel}

% --- Wissenschaftliche Pakete ---
\usepackage{cite}
\usepackage{url}
\usepackage{booktabs}

\begin{document}

\title{Untersuchung mikroarchitektonischer Ressourcen-Interferenzen (Noisy-Neighbor) in virtualisierten Umgebungen}

\author{%
  \IEEEauthorblockN{Maximilian Wagner}
  \IEEEauthorblockA{%
    \textit{Studiengang IT-Sicherheit} \\
    \textit{Fachbereich Informatik und Medien} \\
    \textit{Technische Hochschule Brandenburg} \\
    maximilian.wagner@th-brandenburg.de
  }
}

\maketitle

\begin{abstract}
Dieses Dokument enthält die überarbeitete Version der Einleitung der
wissenschaftlichen Hausarbeit zum Thema mikroarchitektonischer
Ressourcen-Interferenzen in virtualisierten Umgebungen. Die Überarbeitung
basiert auf einer systematischen wissenschaftlichen Analyse der ersten
Rohfassung. Im Anhang ist das Überarbeitungs-Protokoll mit den
vorgenommenen strukturellen und inhaltlichen Korrekturen dokumentiert.
\end{abstract}

% =============================================================================
\section{Einleitung}
% =============================================================================

Virtualisierung bezeichnet die softwaregestützte Abstraktion physischer
Rechenhardware in logisch isolierte Ausführungsumgebungen, sogenannte
Gastsysteme. Ein Hypervisor vermittelt dabei zwischen den Gastsystemen
und der physischen Hardware, um zu gewährleisten, dass Instanzen weder
den Adressraum noch die Ausführungskontrolle benachbarter Systeme direkt
adressieren können. Diese logische Isolation bildet das tragende
Sicherheitsprinzip von Multi-Tenant-Architekturen in modernen
Cloud-Infrastrukturen. Auf mikroarchitektonischer Ebene bleibt diese
Isolation jedoch unvollständig: Physische Komponenten wie der Last Level
Cache (LLC), die Hauptspeicherbandbreite und I/O-Subsysteme werden von
koexistierenden Systemen geteilt. Da der Hypervisor diese
Hardwarekomponenten oft transparent an die Gastsysteme weiterreicht,
entstehen unkontrollierte Interferenzkanäle \cite{koh2007analysis}.

Diese unkontrollierte Ressourcenteilung begründet eine systemische
Schwachstellenklasse mit erheblichem Gefährdungspotenzial. Während
Virtualisierungslösungen eine funktionale Fehler- und Sicherheitsisolation
bereitstellen, kann eine strikte Performance-Isolation strukturell oft nicht
garantiert werden \cite{koh2007analysis}. Die resultierende
Ressourcen-Konkurrenz manifestiert sich insbesondere im sogenannten
Noisy-Neighbor-Effekt: Ein aggressiver Workload innerhalb einer Instanz
verdrängt gezielt Cache-Zeilen koexistierender Systeme. Dies führt zu einer
unautorisierten Degradation von Durchsatz und Latenz benachbarter Mandanten,
was einem mikroarchitektonischen Denial-of-Service-Angriff gleichkommt
\cite{ge2018survey}. Regulatorische Standards wie die Richtlinien des
National Institute of Standards and Technology (NIST) klassifizieren diese
unkontrollierte Ressourcenerschöpfung auf Plattformebene daher als kritisches
Risiko und fordern den Einsatz hardwaregestützter Isolations- und
Kontrollmechanismen \cite{nist2014hypervisor}.

Während die grundlegende Anfälligkeit geteilter Ressourcen in der Literatur
dokumentiert ist, besteht bei aktuellen Hardware-Architekturen ein Defizit
an systematischen Untersuchungen. Moderne Prozessorgenerationen
implementieren komplexe, hybride Kern-Topologien (Performance- und
Efficient-Cores) sowie hardwaregestützte Abwehrfunktionen wie die Intel
Resource Director Technology (RDT) inklusive Cache Allocation Technology
(CAT). Inwiefern diese Schutzmechanismen unter produktiven Bedingungen eine
hinreichende Isolation gewährleisten, ist bisher unzureichend dokumentiert.
Die vorliegende Arbeit adressiert diese Lücke und prüft die Hypothese, ob
historische, LLC-basierte Interferenz-Vektoren auf einer aktuellen Intel
Raptor Lake Architektur unter Verwendung eines produktiven Hypervisors
(Proxmox VE) zu statistisch quantifizierbaren Isolationsbrüchen führen.

Der zur Überprüfung gewählte Messansatz basiert auf einem kontrollierten
Verfahren auf Host-Ebene. Um Verfälschungen der Ergebnisse durch dynamische
Hardware-Optimierungen auszuschließen, werden auf dem Host-System
Energiesparmodi (C-States) sowie die dynamische Frequenzskalierung (Intel
Turbo Boost, CPU-Governor) vollständig deaktiviert. Auf diesem bereinigten
Fundament baut die Untersuchung auf: Die methodische Basis bildet die
Evaluierung und der direkte Leistungsvergleich der fundamentalen
Virtualisierungsparadigmen – der reinen Emulation, der Para-Virtualisierung
und der hardwareunterstützten Virtualisierung – hinsichtlich ihrer Rechen-
und I/O-Kapazitäten.

Auf diesem Fundament wird der spezifische Noisy-Neighbor-Proof-of-Concept
(PoC) aufgesetzt. Hierbei werden zwei Debian-Instanzen (Angreifer und
Opfer) mittels CPU-Affinity auf identische physische Performance-Kerne
gebunden, um eine Cache-Konkurrenz zu erzwingen. Während das
Angreifer-System über synthetische Lasten eine Sättigung des LLC provoziert,
erfasst das Opfer-System mittels standardisierter Benchmarking-Werkzeuge die
Veränderungen von Durchsatz und Latenz. Sollten die hardwareseitigen
Schutzmechanismen des Prozessors (Intel RDT/CAT) die Interferenzen
vollständig neutralisieren und die Hypothese somit falsifiziert werden, greift
eine strukturierte Fallback-Strategie: In diesem Szenario wird die Arbeit auf
den umfassenden Leistungsvergleich der drei Virtualisierungsparadigmen (QEMU,
LXC und KVM) unter standardisierten Rechen- und I/O-Lasten fokussiert, um die
grundlegenden Isolationseigenschaften und Performance-Overheads systematisch
zu quantifizieren.

Die weitere Arbeit gliedert sich wie folgt: Abschnitt II bereitet den aktuellen
Stand der Forschung zu mikroarchitektonischen Interferenz-Vektoren und
Virtualisierungskonzepten auf. Abschnitt III dokumentiert detailliert den
Versuchsaufbau sowie die Host- und Gast-Konfigurationen. Die Präsentation und
kritische Evaluation der Messergebnisse des PoC sowie der Vergleiche erfolgt
in Abschnitt IV. In Abschnitt V werden die Implikationen der Ergebnisse im
Kontext der Ausgangshypothese diskutiert, bevor Abschnitt VI die Arbeit mit
einer Zusammenfassung und einem Ausblick auf zukünftige Forschungsfragen
abschließt.

\bibliographystyle{IEEEtran}
\bibliography{references}

% =============================================================================
% Ab hier: Anhang in Einzelspalte
% =============================================================================
\clearpage
\onecolumn

\section*{Anhang: Erstellungs-Dokumentation}

\noindent
Dieser Anhang dokumentiert den Überarbeitungsprozess der zweiten Version
auf Basis einer systematischen qualitativen Analyse der ersten Rohfassung.

% --- Protokoll -----------------------------------------------------------------
\subsection*{A. Protokoll der wissenschaftlichen Analyse}

\noindent
\begin{tabular}{@{}ll@{}}
  \textbf{Untersuchungsgegenstand:} &
    Überarbeitete Fassung der Einleitung auf Basis der Analyse der Rohfassung \\
  \textbf{Referenzdokument:}        &
    Protokoll der Analyse (Erste Version) \\
  \textbf{Methode:}                 &
    Systematische qualitative Überarbeitung nach IEEE-Kriterien \\
\end{tabular}

\bigskip
\noindent\textbf{A.1\quad Wiederherstellung der inhaltlichen Validität und
Objektivität}

\begin{itemize}
  \item \textbf{Maßnahme:} Der Titel wurde auf den präzisen Wortlaut des
        Exposés (,,Untersuchung mikroarchitektonischer
        Ressourcen-Interferenzen (Noisy-Neighbor) in virtualisierten
        Umgebungen``) zurückgesetzt, um den Untersuchungsgegenstand korrekt
        zu repräsentieren.
  \item \textbf{Begründung:} Der vom LLM generierte Titel verschob den Fokus
        inkorrekt in Richtung Side-Channels und wirkte reißerisch ---
        konträr zum Exposé.
  \item \textbf{Ergebnis:} Der sprachliche Stil versachlicht und strikt
        auf akademische Präzision (Ökonomie des Ausdrucks) ausgerichtet,
        um den Lesefluss zu verbessern.
\end{itemize}

\bigskip
\noindent\textbf{A.2\quad Korrektur des Literaturfokus
(Forschungstransparenz)}

\begin{itemize}
  \item \textbf{Maßnahme:} Die namentliche, argumentative
        Auseinandersetzung mit spezifischen Autoren im Fließtext wurde
        entfernt.
  \item \textbf{Begründung:} Eine detaillierte Literatursynthese ist
        methodisch dem Kapitel ,,Stand der Forschung`` vorbehalten; die
        Einleitung darf diese nicht vorwegnehmen.
  \item \textbf{Ergebnis:} Problemstellung abstrahiert; Quellen dienen nun
        ausschließlich als komprimierter Nachweis in Klammern.
\end{itemize}

\bigskip
\noindent\textbf{A.3\quad Anpassung des Abstraktionsniveaus}

\begin{itemize}
  \item \textbf{Maßnahme:} Technologische Detailtiefe reduziert
        (z.\,B.\ Streichung spezifischer Angriffsvektoren wie
        PRIME+PROBE).
  \item \textbf{Begründung:} Eine Einleitung operiert auf konzeptioneller
        Ebene; die Nennung konkreter Angriffsmethoden greift der späteren
        Detailanalyse unzulässig vor.
  \item \textbf{Ergebnis:} Fokus liegt wieder auf der konzeptionellen
        Abgrenzung des Problemfeldes der logischen Isolation gemäß Exposé.
\end{itemize}

\bigskip
\noindent\textbf{A.4\quad Strukturelle und methodische Kohärenz}

\begin{itemize}
  \item \textbf{Maßnahme:} Die starre Abarbeitung expliziter Leitfragen
        (Aufzählungsartefakte) wurde in einen organischen, zusammenhängenden
        Fließtext überführt.
  \item \textbf{Begründung:} IEEE-konforme Einleitungen erfordern einen
        kohärenten Prosafluss ohne explizite Zwischenüberschriften.
  \item \textbf{Ergebnis:} Forschungslogik zur Fallback-Strategie
        geschärft --- es wird nun methodisch sauber hergeleitet, dass der
        allgemeine Leistungsvergleich der Virtualisierungslösungen
        (Basisaufgabe) das fundierte Messbett bildet, auf dem der
        spezifische Noisy-Neighbor-Angriff anschließend evaluiert wird.
\end{itemize}

\end{document}